\bookchapter{บทที่ 22: AI ทำงานอย่างไรในโทรศัพท์มือถือ}

\chapterintro{ยินดีต้อนรับสู่ “เวทีมายากล” สุดเจ๋งของมือถือ โดยมีพระเอกคือปัญญาประดิษฐ์ (AI) ไม่ว่าจะเป็นภาพถ่ายสวยที่กดถ่ายปุ๊บก็ได้ปั๊บ ผู้ช่วยเสียงที่พูดปุ๊บตอบปั๊บ หรือแป้นพิมพ์ที่ยิ่งใช้ยิ่งเข้าใจคุณ ทั้งหมดนี้ล้วนมี “สมองล่องหน” ตัวนี้อยู่เบื้องหลัง}

\booksection{หนึ่ง กล้อง AI: “ห้องมืดตกแต่งภาพ” ที่ย้ายเข้ามาอยู่ในมือถือของคุณแล้ว}

การถ่ายภาพด้วย AI ยุคนี้ เทียบได้กับการยัดทีมช่างภาพมืออาชีพบวกช่างตกแต่งภาพเข้าไปไว้ในปุ่มชัตเตอร์เลยทีเดียว ก่อนที่คุณจะกดชัตเตอร์ AI เริ่มทำงานไปแล้ว: สแกนฉาก จำแนกออกมาว่า “นี่คือท้องฟ้าสีคราม” “นี่คือใบหน้าคน” “นี่คือภาพมุมกลางคืน”

\booksection{สอง ผู้ช่วยเสียง: “เลขานุการประจำเสียง” ของคุณเข้าใจคุณได้อย่างไร}

\begin{enumerate}
  \item \textbf{ได้ยินชัด (การจดจำเสียงพูด)}: เปลี่ยนคลื่นเสียงให้กลายเป็นข้อความตัวอักษร
  \item \textbf{เข้าใจความหมาย (การเข้าใจภาษาธรรมชาติ)}: เข้าใจเจตนา และดึงสารสำคัญออกมา
  \item \textbf{ลงมือทำและตอบ}: เรียกใช้แอปหรือบริการที่เกี่ยวข้อง แล้วสังเคราะห์เสียง "พูด" บอกคุณอีกที
\end{enumerate}

\booksection{สาม AI บนเครื่อง vs. AI บนคลาวด์}

AI บนเครื่องเก็บการประมวลผลไว้ในมือถือ ทั้งเร็ว ปลอดภัยต่อความเป็นส่วนตัว และใช้แบบออฟไลน์ได้ ส่วน AI บนคลาวด์ต้องต่อเน็ต แต่รับมือกับงานซับซ้อนกว่าได้

\booksection{สี่ โมเดล AI ขนาดใหญ่กับมือถือ: ผู้ช่วยอัจฉริยะประจำตัวของคุณ}

ช่วงปี 2025-2026 การระเบิดของโมเดล AI ขนาดใหญ่ (อย่าง DeepSeek R1, ChatGPT-5, ERNIE Bot 5.0, Qwen 3 ฯลฯ) ได้เปลี่ยนวิธีใช้งานมือถือไปอย่างลึกซึ้ง มือถือไม่ใช่แค่เครื่องมือรับคำสั่งอีกต่อไป แต่กลายเป็นคู่หูอัจฉริยะที่เข้าใจ ใช้เหตุผล และสร้างสรรค์ได้ ชิป AI บนเครื่อง (พลัง NPU ทะลุ 80 TOPS) ทำให้การรันโมเดลใหญ่บนตัวเครื่องเองเป็นจริงได้เป็นครั้งแรก ไม่ต้องต่อเน็ตก็ใช้ AI ได้

\booksubsection{โมเดล AI ขนาดใหญ่คืออะไร?}

โมเดล AI ขนาดใหญ่ (Large Language Model: LLM) คือระบบ AI ที่ฝึกมาจากข้อมูลข้อความมหาศาล ทำความเข้าใจและสร้างภาษาธรรมชาติได้ ลองนึกว่ามันคือ “ตัวเก่งพิเศษ” ที่อ่านข้อความบนอินเทอร์เน็ตมาแทบทุกชิ้น ตอบคำถาม เขียนบทความ แปลภาษา เขียนโปรแกรม วิเคราะห์ข้อมูลได้หมด

\booksubsection{การใช้งานโมเดล AI ใหญ่บนมือถือ:}

\begin{itemize}
  \item \textbf{ผู้ช่วยสนทนา}: แอปอย่าง DeepSeek, ChatGPT, Doubao, Kimi ให้ถามได้ทุกเมื่อ แปลภาษา เขียนคอนเทนต์ สรุปบทความ ฉลาดกว่าเสิร์ชเอนจินแบบเดิม เพราะไม่ได้แค่โยนลิงก์ให้คุณ แต่ตอบคำตอบตรง ๆ เลย
  \item \textbf{การเขียนและงานสร้างสรรค์}: ให้ AI ช่วยร่างอีเมล สรุปมโนทัศน์ที่ประชุม ขัดเกลาบทความ คิดคอนเทนต์เก๋ ๆ
  \item \textbf{สร้างและแต่งรูป}: AI ลบพวกคนแปลกหน้าในภาพออกแบบเสร็จเดียว ขยายขอบภาพ เปลี่ยนสไตล์ กระทั่งสร้างรูปจากคำบรรยายได้
  \item \textbf{แปลภาษาแบบเรียลไทม์}: แปลระหว่างสายสด ๆ แปลบทสนทนาเจอหน้ากัน แปลจากป้ายที่ถ่ายรูป เที่ยวต่างประเทศไม่มีกำแพงภาษาอีกต่อไป
  \item \textbf{ติวเตอร์ส่วนตัว}: AI อธิบายคอนเซปต์ ออกโจทย์ให้ฝึก ตรวจการบ้าน เป็นได้ทั้งติวเตอร์ที่คอยอยู่กับคุณ
\end{itemize}

\booksubsection{โมเดลใหญ่บนเครื่อง: มือถือรัน AI เองได้}

มือถือรุ่นเรือธงล่าสุดรันโมเดล AI ขนาดย่อม (โดยทั่วไปราว 7B พารามิเตอร์) ในเครื่องได้แล้ว สิ่งที่เปลี่ยนไปคือ:

\begin{itemize}
  \item \textbf{ไม่ต้องพึ่งเน็ต}: เน็ตหลุดก็ยังใช้ฟีเจอร์ AI ได้ ทั้งบนเครื่องบินและในรถไฟใต้ดิน
  \item \textbf{ความเป็นส่วนตัวปลอดภัยกว่า}: บทสนทนาและข้อมูลของคุณไม่ถูกอัปโหลดขึ้นคลาวด์ ประมวลผลอยู่ในเครื่องล้วน ๆ
  \item \textbf{ตอบไวกว่า}: ไม่มีดีเลย์จากเครือข่าย AI ตอบแทบจะทันทีทันควัน
  \item \textbf{ข้อจำกัดในปัจจุบัน}: ความสามารถของโมเดลบนเครื่องยังสู้โมเดลใหญ่บนคลาวด์ไม่ได้ งานซับซ้อนก็ยังต้องพึ่งเน็ต แต่เมื่อพลัง AI ของชิปแรงขึ้นเรื่อย ๆ โมเดลบนเครื่องก็จะเก่งขึ้นเรื่อย ๆ เหมือนกัน
\end{itemize}

\booksubsection{การผสาน AI ระดับระบบปฏิบัติการ:}

\begin{itemize}
  \item \textbf{Apple Intelligence ของแอปเปิล}: ระบบ AI ที่มากับ iOS 18 รวมเครื่องมือช่วยเขียน ไม้กายสิทธิ์แต่งรูป Siri ฉบับเสริมพลัง สรุปการแจ้งเตือน ฯลฯ ให้ความสำคัญกับการประมวลผลบนเครื่องเป็นหลัก งานซับซ้อนถึงจะส่งขึ้นคลาวด์
  \item \textbf{โมเดล Pangu ของหัวเหวย}: ผสานเข้ากับ HarmonyOS ผู้ช่วย Xiao Yi ได้พลังระดับโมเดลใหญ่ ทั้งสรุปเอกสาร อ่านหน้าจออัจฉริยะ แต่งรูปด้วย AI ฯลฯ
  \item \textbf{HyperAI ของ Xiaomi}: ความสามารถ AI ระดับระบบ ทั้งค้นรูปด้วย AI สรุปด้วย AI คำบรรยายสดด้วย AI และแชทกับโมเดลใหญ่บนเครื่อง
  \item \textbf{Galaxy AI ของซัมซุง}: แปลระหว่างสายสด ๆ สรุปโน้ตด้วย AI สร้างและแต่งรูป วงค้นหา (Circle to Search) ฯลฯ
\end{itemize}

\booksection{ห้า ตัวอย่างสถานการณ์ใช้งานจริงและการลงมือทำ}

เพื่อให้คุณเห็น “พลังพิเศษ” ของ AI ในมือถืออย่างเห็นภาพ เราขอยกห้าสถานการณ์จริงต่อไปนี้ มาโชว์ว่า AI “เข้าช่วยเหลือ” ชีวิตประจำวันได้อย่างไร:

\booksubsection{ฉากที่ 1: แปลภาษาเรียลไทม์ — ล่ามประจำตัวสำหรับเที่ยวต่างประเทศ}

\textbf{การใช้งานจริง}: สั่งอาหารที่ปารีสแต่พูดฝรั่งเศสไม่ได้? แอป “แปลภาษา” ของ iPhone มีโหมดสนทนาตัวต่อตัว แต่ละฝ่ายอ่านคำบรรยายเป็นภาษาของตัวเอง แถมแปลสดต่อเนื่องได้ด้วย ฝั่งแอนดรอยด์อย่าง Galaxy AI ของซัมซุงก็แปลระหว่างสายได้ คุณพูดภาษาจีนเข้าไป อีกฝ่ายได้ยินเป็นภาษาอังกฤษ และกันกลับทางก็ได้เช่นกัน

\textbf{หลักการเบื้องหลัง}: ไปป์ไลน์ 3 ขั้น คือ การจดจำเสียงพูด (ASR) การแปลด้วยเครื่อง (NMT) การสังเคราะห์เสียง (TTS) ความหน่วงตั้งแต่ต้นจนจบถูกลดเหลือภายใน 1 วินาทีได้แล้ว

\booksubsection{ฉากที่ 2: AI ลบวัตถุ — ช่วยชีวิตภาพถ่ายเที่ยวที่โดนคนแปลกหน้า “ทำลาย”}

\textbf{การใช้งานจริง}: ถ่ายเซลฟี่สวยสุด ๆ ที่จุดเช็กอิน แต่ฉากหลังแน่นไปด้วยนักท่องเที่ยวด้วยกันเอง? ใช้ฟีเจอร์ “ลบแบบมายากล” วงคนแปลกหน้าไว้ มือถือจะจำแนกแล้ว “ป้าย” คนเหล่านั้นออก จากนั้นเติมฉากหลังให้เนียนอัตโนมัติ ผลลัพธ์มักน่าทึ่งจนต้องแปลกใจ

\textbf{หลักการเบื้องหลัง}: โมเดลแบ่งภาพและเติมภาพ (Inpainting) จะหาพื้นที่ที่เลือกแล้วลบทิ้งก่อน จากนั้นอาศัยพิกเซลรอบข้างและความเข้าใจเชิงความหมาย “นึกภาพ” สิ่งที่ควรเติมขึ้นมาแทน โมเดล AI เบื้องหลังเทคนิคนี้ฝึกด้วยรูปมากถึงหลายล้านใบ

\booksubsection{ฉากที่ 3: แต่งรูปด้วย AI — แตะเดียวกลายเป็น “ช่างภาพมือรางวัล”}

\textbf{การใช้งานจริง}: ได้ภาพวิววันเมฆหม่น ๆ มาใบหนึ่ง? ฟีเจอร์ “ท้องฟ้ามายากล” ของ AI เปลี่ยนท้องฟ้าเป็นแดดใส สีพระอาทิตย์อัสดง หรือแม้ท้องฟ้าดาว ได้ในคลิกเดียว ส่วน AI ขยายภาพช่วยเปลี่ยนภาพแนวตั้งเป็นแนวนอนได้ โดยมือถือ “นึกภาพ” สิ่งที่อยู่นอกกรอบให้เอง

\textbf{หลักการเบื้องหลัง}: ใช้โมเดลการแพร่ (Diffusion Model) จาก AI เชิงสร้างสรรค์ (Generative AI) เข้าใจเนื้อหาเชิงความหมายของภาพแล้วสร้างพิกเซลใหม่ขึ้นมา จัดเป็น “แนวหน้า” ของงานประมวลผลภาพด้วย AI ในยุคนี้

\booksubsection{ฉากที่ 4: จัดการเอกสาร — เครื่องคูณประสิทธิภาพของคนออฟฟิศ}

\textbf{การใช้งานจริง}: ประชุมเสร็จ โยนไฟล์บันทึกเสียงทิ้งไว้ให้ AI มันจะจัดระเบียบเป็นมโนทัศน์ที่ประชุมแบบมีโครงเรื่อง ดึงสิ่งที่ต้องทำออกมา พร้อมสรุปย่อให้ รายงาน PDF ภาษาอังกฤษ 10 หน้า AI คัดข้อสรุปสำคัญให้ภายใน 10 วินาที

\textbf{หลักการเบื้องหลัง}: โมเดลบริบทยาว (Long Context LLM) “อ่าน” เอกสารหลายหมื่นคำจบได้ แล้วใช้เหตุผลหลายรอบพร้อมดึงสารสำคัญ Apple Intelligence, Xiao Yi ของหัวเหวย, WPS AI, Tencent Yuanbao ฯลฯ ล้วนมีความสามารถแบบนี้

\booksubsection{ฉากที่ 5: ค้นรูปและจำแนกสิ่งของด้วย AI — “เสิร์ชเอนจินสายตา” ส่วนตัวของคุณ}

\textbf{การใช้งานจริง}: เห็นดอกไม้แปลก ๆ ริมถนน เปิดกล้องขึ้นมา AI บอกได้ทันทีว่า “นี่คือดอกชบา” เวลาไล่หารูปที่ถ่ายไว้เมื่อปีกลางริมทะเล พิมพ์ค้นว่า “ทะเล” AI กวาดรูปที่มีทะเลทั้งหมดขึ้นมาให้เอง ไม่ต้องเลื่อนหาเองทีละใบ

\textbf{หลักการเบื้องหลัง}: โมเดลคอมพิวเตอร์วิทัศน์ (CV) แปลงเนื้อหาในรูปเป็นเวกเตอร์ (Embedding) แล้วจับคู่ความคล้ายกับคำค้น การ “ค้นหาแบบหลายมิติ” นี้คือก้าวกระโดดครั้งใหญ่ของการจัดการอัลบั้มบนมือถือ

\booksection{หก ปรับแต่งพรอมป์: “สูตรสื่อสาร” ที่ทำให้ AI ทำตามคำสั่งของคุณ}

คุณภาพคำตอบของ AI ขึ้นอยู่กับว่าคุณ “ถาม” อย่างไรเป็นหลัก คำถามเดียวกัน ถามคนละแบบ คำตอบก็ต่างกันลิบลับ ต่อไปนี้เราจะแนะนำ\textbf{แนวคิดสื่อสารแบบ “บทบาท-พื้นหลัง-งาน-ข้อจำกัด” ที่ใช้ได้จริงและเริ่มง่ายที่สุด}พร้อมยกตัวอย่างประกอบมากมาย

\begin{tcbinfo}{ทำไมต้องมีกรอบ?}
คนจำนวนมากใช้ AI เหมือนใช้ Baidu — พิมพ์คำสำคัญไม่กี่คำแล้วรอผล แต่ AI ไม่ใช่เสิร์ชเอนจิน มันคือโมเดลภาษาที่ “อ่านจับใจความแล้วจึงตอบ”\textbf{ให้บริบทมากขึ้น บอกให้ชัดว่าต้องการอะไร และกำหนดขอบเขตไม่ให้เพี้ยน}เมื่ออย่างนั้นมันจะให้ผลลัพธ์ที่แม่นยำ กรอบพรอมป์ก็คือการแปล “คำตอบที่ดีในหัวของคุณ” ให้กลายเป็นคำสั่งที่ AI เข้าใจ

\end{tcbinfo}

\booksubsection{บทบาท (Role): บอก AI ว่ามันคือใคร}

เมื่อให้ “ตัวตน” ที่เฉพาะเจาะจงแก่ AI มันจะดึงความรู้และวิธีสื่อสารของสายนั้นมาใช้\textbf{ยิ่งตั้งบทบาทเฉพาะเจาะจงเท่าไร ผลลัพธ์ยิ่งเป็นมืออาชีพเท่านั้น}

\booksubsection{ตัวอย่างใช้ได้จริง:}

\begin{itemize}
  \item "คุณคือ\textbf{วิศวกรซ่อมมือถือ Xiaomi ประสบการณ์ 10 ปี}." (คำถามสายเทคนิค)
  \item "คุณคือ\textbf{ครูฟิสิกส์มัธยมปลายผู้ชำนาญอธิบายเรื่องยากให้ง่าย}." (คำถามสายวิทยาศาสตร์)
  \item "คุณคือ\textbf{บรรณาธิการหนังสือสายตาดี อ่านกว้าง มีสไตล์เป็นของตัวเอง}." (คำถามสายการเขียน)
  \item "คุณคือ\textbf{บล็อกเกอร์อาหารที่ทำเมนูมื้อเช้าให้คนทำงานทุกวัน}." (คำถามสายชีวิตประจำวัน)
  \item "คุณคือ\textbf{หมอแพทย์แผนจีนผู้ใจเย็น อ่อนโยน และไม่เคยหัวเราะเยาะมือใหม่}." (คำถามสายสุขภาพ)
\end{itemize}

\textbf{ทำไมจึงได้ผล?}การตั้งบทบาทช่วยกำหนด\textbf{ขอบเขตความรู้}\textbf{สไตล์การพูด}และ\textbf{ระดับความเชี่ยวชาญ}ให้แก่ AI ส่วน AI ที่ไม่มีบทบาทเป็นแค่ “นักเรียนที่รู้กว้างแต่ไม่เชี่ยวชาญอะไรลึก” พอมีบทบาท มันจะ “เข้าบท” สวมบทเป็นมืออาชีพสายนั้น แล้วคิดเหมือนเขาจริง ๆ

\booksubsection{พื้นหลัง (Background): เล่าสถานการณ์ของคุณให้ AI ฟัง}

บอกให้ชัดตั้งแต่ต้นว่า “คุณคือใคร ตอนนี้สถานการณ์เป็นอย่างไร ติดข้อจำกัดอะไร” AI จะให้คำตอบ\textbf{ที่เหมาะกับคุณจริง ๆ}ไม่ใช่คำตอบลอย ๆ ก้ำกึ่ง

\booksubsection{ตัวอย่างใช้ได้จริง:}

\begin{itemize}
  \item "มือถือของฉันคือ\textbf{Xiaomi 14 ซื้อมาได้ครึ่งปี ช่วงนี้เล่นเกมแล้วเครื่องร้อนมาก}." (ข้อมูลเครื่อง)
  \item "ฉันเป็น\textbf{คุณป้าวัย 60 ที่เกษียณแล้ว สายตาไม่ค่อยดี ใช้สมาร์ตโฟนไม่ค่อยเป็น}." (ภาพผู้ใช้)
  \item "ฉันกำลัง\textbf{เตรียมสอบวัดระดับภาษาอังกฤษ CET-6 เดือนมิถุนายน 2026 พื้นฐานก็เรื่อย ๆ}." (พื้นหลังการเรียน)
  \item "ฉันเปิด\textbf{ร้านกาแฟเล็ก ๆ ในเมืองรอง รายได้เดือนละ 30,000 มีลูกจ้างสองคน}." (พื้นหลังธุรกิจ)
  \item "ฉัน\textbf{เพิ่งคลอดลูก อยู่บ้านพักฟื้น นอนไม่พอหนักมาก}." (ช่วงชีวิต)
\end{itemize}

\textbf{ทำไมจึงได้ผล?}ข้อมูลพื้นหลังทำให้ AI รู้ว่าต้อง\textbf{“รับใช้ใคร”}และ\textbf{“อยู่ภายใต้เงื่อนไขอะไร”}คำถามเดียวกัน การอธิบายให้นักเรียนกับให้ผู้ปกครองย่อมคนละแบบ คำแนะนำให้คนเมืองกับคนต่างจังหวัดก็ต่างกัน ถ้าไม่มีพื้นหลัง AI จะได้แต่คำตอบ “สามัญที่สุด” ซึ่งมักจืดชืดไม่เต็มที่

\booksubsection{งาน (Task): บอกให้ชัดว่าให้ AI ทำอะไร}

ขั้นนี้ต้อง\textbf{เริ่มด้วยกริยา}(เขียน จัดรายการ เปรียบเทียบ แปล อธิบาย สรุป...) และอธิบายงานให้\textbf{เจาะจง ลงมือทำได้จริง}งานที่กำกวม ย่อมได้คำตอบกำกวม ๆ

\booksubsection{ตัวอย่างใช้ได้จริง:}

\begin{itemize}
  \item "กรุณา\textbf{จัดรายการ 5 สิทธิ์ที่ฉันควรปิดทันที}พร้อมบอกเหตุผล." (การกระทำชัดเจน)
  \item "กรุณา\textbf{แปลอีเมลภาษาอังกฤษข้างล่างนี้เป็นภาษาไทย}รักษาโทนธุรกิจ." (งานเจาะจง)
  \item "กรุณา\textbf{เปรียบเทียบมือถือ 4 รุ่นในช่วงราคา 3,000 หยวนนี้}พร้อมทำตารางข้อดีข้อเสียให้." (รูปแบบชัดเจน)
  \item "กรุณา\textbf{อธิบายว่า 5G คืออะไรด้วยภาษาที่เด็กประถมก็เข้าใจ}." (สไตล์ชัดเจน)
  \item "กรุณา\textbf{เขียนคำบรรยายคลิปสั้น 3 ชิ้นให้ร้านมื้อเช้าของฉัน}โดยเจาะแยกกลุ่มคนทำงาน คู่รัก และพ่อแม่พาลูก." (การแบ่งกลุ่มชัดเจน)
\end{itemize}

\textbf{ทำไมจึงได้ผล?}ส่วนของงานคือ “ศูนย์บัญชาการ” ของการลงมือของ AI\textbf{กริยายิ่งชัด ข้อกำหนดยิ่งเจาะจง AI ยิ่งรู้ว่า “ต้องส่งมอบอะไร”}คำหลวม ๆ อย่าง “ช่วยจัดการหน่อย” “ทำสักอย่าง” “ลองดูสิ” เป็นสิ่งต้องเลี่ยงสุด ๆ เมื่อใช้ AI มันจะเดาใจคุณ แล้วผลลัพธ์มักเพี้ยน

\booksubsection{ข้อจำกัด (Constraints): บอก AI ว่าไม่ต้องทำอะไร}

“ข้อจำกัด” ฟังดูเหมือนการมัดมือมัดเท้า แต่จริง ๆ คือเครื่องมือ\textbf{บีบ AI ให้ลงล็อกที่ผลลัพธ์ที่เหมาะกับคุณที่สุด}มิติที่ใช้บ่อยมี: จำนวนคำ รูปแบบ น้ำเสียง สิ่งที่เลี่ยง และสิ่งที่ต้องมี

\booksubsection{ตัวอย่างใช้ได้จริง:}

\begin{itemize}
  \item "\textbf{ไม่เกิน 200 คำ จัดเป็นหัวข้อย่อย}." (ความยาว+รูปแบบ)
  \item "\textbf{น้ำเสียงต้องอบอุ่นเป็นกันเอง เหมือนเพื่อนคุยกัน}ห้ามใช้ภาษาทางการแบบ "ประการแรก ประการที่สอง ประการที่สาม"." (น้ำเสียง)
  \item "\textbf{ห้ามใช้ศัพท์เทคนิค}ถ้าเจอคำเฉพาะทางให้อธิบายผ่านเลย." (ข้อจำกัดคำศัพท์)
  \item "\textbf{แต่ละย่อหน้าไม่เกิน 3 บรรทัด}เพื่ออ่านสบายบนมือถือ." (การจัดหน้า)
  \item "\textbf{ต้องมี: สถานการณ์การใช้งาน ข้อดีข้อเสีย และความคุ้มกับงบของฉัน}." (สิ่งที่ต้องมี)
  \item "\textbf{ห้ามแนะนำการลงทุนเฉพาะเจาะจง}อธิบายแค่หลักการพอ." (เส้นแบ่งขอบเขต)
\end{itemize}

\textbf{ทำไมจึงได้ผล?}ไม่มีข้อจำกัด AI มัก “แสดงฝีมือเกินงาน” — อย่างเป็นวิชาการเกินไป ยาวเกินไป หรือออกนอกเรื่อง ข้อจำกัดที่ชัดเจน = เกณฑ์การส่งมอบที่ชัดเจน ทั้งสองฝ่ายตรงกันที่ความคาดหวัง ผลลัพธ์จึงควบคุมได้

\booksubsection{ประเมินคุณภาพพรอมป์: 5 ลักษณะของพรอมป์ที่ดี}

เขียนพรอมป์ให้เนี้ยบไม่ใช่เรื่องเขียนครั้งเดียวจบ\textbf{พอเขียนเสร็จ ให้เช็กห้าข้อนี้ก่อน}แล้วค่อยส่งให้ AI:

\begin{enumerate}
  \item \textbf{เจาะจง ไม่กำกวม}: ไม่ใช่ "เขียนบทความหน่อย" แต่คือ "เขียนคู่มือเลือกซื้อมือถือสำหรับแม่ลูกอ่อน 800 คำ แบ่ง 3 ตอนย่อย"
  \item \textbf{มีบริบท ไม่ลอย ๆ}: บอก AI ว่า "ฉันคือใคร ทำไมถึงถามเรื่องนี้" มันจึงให้คำตอบแบบ "ตัดเย็บพอดีตัว"
  \item \textbf{รูปแบบผลลัพธ์ชัดเจน}: จะเอาเป็นตาราง รายการข้อ ๆ ย่อหน้า หรือ JSON? บอกไว้ก่อนเลย
  \item \textbf{ข้อจำกัดครบถ้วน}: ความยาว สไตล์ สิ่งที่เลี่ยง ใส่แค่ข้อใดข้อหนึ่งก็เห็นความต่างชัดเจนแล้ว
  \item \textbf{ตรวจสอบกับ AI ได้}: ถามเสร็จ ให้ AI พูดย้ำว่าเข้าใจความต้องการของคุณอย่างไร ถ้ามันเล่าไม่ชัด แปลว่าคุณเองก็ยังบอกไม่ชัด
\end{enumerate}

\booksubsection{ปัญหาที่พบบ่อยและวิธีแก้}

\begin{center}
\begin{tabularx}{\linewidth}{|X|X|}
\hline
\textbf{อาการ} & \textbf{วิธีแก้} \tabularnewline
\hline
คำตอบกว้างเกินไป ใช้ไม่ได้จริง & เพิ่มพื้นหลัง (คุณคือใคร/อยู่ในสถานการณ์แบบไหน) \tabularnewline
คำตอบเป็นวิชาการ/ทางการเกินไป & ใส่ข้อจำกัดน้ำเสียง + สั่งห้ามภาษาทางการแบบราชการ \tabularnewline
คำตอบออกนอกเรื่อง/ตอบไม่ตรงคำถาม & พูดงานซ้ำด้วยกริยาชัดเจน และสั่งว่า "ตอบแค่เรื่องนี้เท่านั้น" \tabularnewline
คำตอบยาวเกินอ่านไม่จบ & กำหนดจำนวนคำแบบตายตัว (ไม่เกิน 200 คำ) \tabularnewline
AI พูดมั่วอย่างมั่นใจเต็มร้อย & ใส่คำสั่ง "ถ้าไม่แน่ใจ ให้ระบุว่า 'ไม่แน่ใจ'" + ตรวจสอบข้อเท็จจริงสำคัญด้วยตัวเอง \tabularnewline
เวอร์ชันแรกไม่ถูกใจ & อย่ารีไรต์ทั้งก้อน ให้ต่อท้ายพรอมป์เดิมด้วย "ทำให้...มากขึ้น" \tabularnewline
\hline
\end{tabularx}
\end{center}

\booksubsection{สามตัวอย่างสถานการณ์ฉบับเต็ม}

\begin{tcbinfo}{สถานการณ์ที่ 1: ผู้ช่วยชีวิตประจำวัน (คำแนะนำจัดเมนูมื้อเช้า)}
\textbf{พรอมป์}:

คุณเป็นบล็อกเกอร์อาหารมื้อเช้าผู้เชี่ยวชาญการจับคู่โภชนาการ (\textbf{บทบาท}) ฉันเป็นแม่ทำงาน ทุกเช้ามีเวลาทำอาหารเช้าแค่ 15 นาที ที่บ้านมีลูกวัย 6 ขวบและแม่สามีวัย 60 ปี (\textbf{พื้นหลัง}) กรุณา\textbf{จัดเมนูอาหารเช้า 7 วันให้ไม่ซ้ำกัน}(\textbf{งาน}) อาหารแต่ละมื้อต้อง: (1) เตรียมไม่เกิน 15 นาที (2) โภชนาการสมดุล (โปรตีน+คาร์บ+ผักผลไม้) (3) ทั้งเด็กและผู้สูงวัยกินได้ (4) หาวัตถุดิบได้จากซูเปอร์มาร์เก็ตทั่วไป (\textbf{ข้อจำกัด})

\textbf{ทำไมจึงได้ผล}: บทบาทกำหนดความเป็นมืออาชีพ พื้นหลังครอบคลุมคนสามช่วงวัย งานระบุสิ่งที่ต้องส่งมอบ ส่วน 4 ข้อจำกัดทำให้ทำตามได้จริง

\end{tcbinfo}

\begin{tcbinfo}{สถานการณ์ที่ 2: ผู้ช่วยเรื่องการเรียน (ติวสอบ CET-6)}
\textbf{พรอมป์}:

คุณเป็น\textbf{ครูสอน CET-6 ที่ประสบการณ์แน่น}(บทบาท) ฉันเป็นนักศึกษาธรรมดาคนหนึ่ง\textbf{CET-4 ได้ 480 คะแนน เป้าหมาย CET-6 อยู่ที่ 500 มีเวลาอีก 3 เดือน}(พื้นหลัง) กรุณา\textbf{วางแผนติว 12 สัปดาห์แบบลงมือทำได้จริง}(งาน) ข้อกำหนด: (1) แบ่งเวลาฟัง พูด อ่าน เขียนรายสัปดาห์ให้ชัดเจน (2) ระบุทิศทางเร่งเครื่องของแต่ละเดือน (3) แนะนำหนังสือทบทวน 3 เล่ม (อธิบายเหตุผลเป็นภาษาไทย) (4) ออกผลเป็นตาราง (ข้อจำกัด)

\textbf{ทำไมจึงได้ผล}: บทบาทกำหนดพื้นเพความเชี่ยวชาญ พื้นหลังปักหมุดคะแนนกับเวลา งานเน้น “แผนเป็นช่วง ๆ” ข้อจำกัดบังคับให้ “เป็นตาราง” และ “ทำได้จริง”

\end{tcbinfo}

\begin{tcbinfo}{สถานการณ์ที่ 3: เพิ่มประสิทธิภาพงานออฟฟิศ (จัดระเบียบมโนทัศน์ที่ประชุม)}
\textbf{พรอมป์}:

คุณเป็น\textbf{ผู้ช่วยบริหารมากประสบการณ์}(บทบาท) ฉันเพิ่งประชุมแผนกยาว 1 ชั่วโมงเสร็จ\textbf{หัวข้อคือแผนรณรงค์การตลาดไตรมาส 2}(พื้นหลัง)\textbf{กรุณาจัดบันทึกการประชุมข้างล่างนี้เป็นมโนทัศน์แบบมีโครงเรื่อง}(งาน) ข้อกำหนด: (1) แบ่งสามส่วนคือ “มติที่ประชุม / สิ่งที่ต้องทำ / ความเสี่ยงและความเห็นที่แตกต่าง” (2) สิ่งที่ต้องทำระบุผู้รับผิดชอบและกำหนดเสร็จด้วย (ถ้าไม่ได้พูดถึง ให้เขียนว่า “รอยืนยัน”) (3) ความยาวรวมไม่เกิน 600 คำ (4) ออกผลเป็น markdown (ข้อจำกัด)

\textbf{ทำไมจึงได้ผล}: บทบาทตั้งมุมมองฝ่ายธุรการ พื้นหลังให้หัวข้อประชุม งานเน้นผลลัพธ์แบบมีโครง 4 ข้อจำกัดทำให้เอาผลไปใช้ต่อได้ทันที

\end{tcbinfo}

\textbf{เทคนิคขั้นสูง: ทำให้ AI เป็น “ผู้ช่วยส่วนตัว” ของคุณ}\par

ลองใช้เทคนิคขั้นสูงเหล่านี้ แล้วจะได้เห็นฝีมือจริงของ AI:

\begin{itemize}
  \item \textbf{ปรับแต่งแบบหลายรอบ}: ไม่พอใจคำตอบของ AI ก็บอกได้ว่า "อันนี้เอาเทคนิคเกินไป ช่วยอธิบายใหม่ด้วยภาษาที่เด็กประถมก็เข้าใจ"
  \item \textbf{การคิดเป็นลูกโซ่ (Chain-of-Thought)}: สั่งให้ AI "คิดทีละขั้น" คุณภาพคำตอบมักดีขึ้น เช่น "ช่วยวิเคราะห์โจทย์คณิตศาสตร์นี้ทีละขั้น"
  \item \textbf{หลาย AI ช่วยกันทำงาน}: ให้ A ร่างแรก B ขัดเกลา C แปลภาษา — ใช้ประกอบกันได้ผลดียิ่งขึ้น
  \item \textbf{คำสั่งระดับระบบ (System Prompt)}: ตั้งค่าตัวตน สิ่งที่ชอบ สิ่งที่เลี่ยง ฯลฯ ไว้ในการตั้งค่าของ AI คำตอบจะเข้ากับความต้องการของคุณยิ่งขึ้น
\end{itemize}

\booksection{เจ็ด ข้อจำกัดและความท้าทายของ AI: รักษาวิจารณญาณให้ตื่นตัว}

AI เก่งจริง แต่ไม่ใช่ “ทำได้ทุกอย่าง” รู้ทันข้อจำกัดของมันอย่างมีสติ จะใช้ได้อย่างเข้าใจมากขึ้น:

\begin{itemize}
  \item \textbf{ปัญหา "ภาพหลอน"}: AI อาจพูดมั่วอย่างจริงจังเต็มร้อย (ศัพท์เทคนิคเรียกว่า Hallucination) มันแต่งเหตุการณ์ประวัติศาสตร์ที่ไม่เคยมีอยู่จริง คำพูดที่ไม่เคยกล่าว หรือตัวเลขสถิติปลอม ๆ ข้อมูลสำคัญต้องตรวจสอบด้วยคนเสมอ
  \item \textbf{ความรู้มีวันตัด}: ข้อมูลฝึกของ AI ส่วนใหญ่มีวันสิ้นสุด (เช่น ต้นปี 2024) เรื่องที่เกิดขึ้นหลังจากนั้นมันไม่รู้
  \item \textbf{ความสามารถใช้เหตุผลจำกัด}: การใช้เหตุผลหลายขั้น การพิสูจน์ทางคณิตศาสตร์ หรือกับดักตรรกะ AI ยังพลาดได้ง่าย
  \item \textbf{ความเสี่ยงด้านความเป็นส่วนตัว}: บริการ AI บนคลาวด์อาจนำบทสนทนาของคุณไปใช้ฝึกโมเดล เรื่องละเอียดอ่อนควรใช้ AI บนเครื่องเป็นหลัก
  \item \textbf{อคติและค่านิยม}: AI อาจมีอคติจากข้อมูลที่ใช้ฝึก ต้องคิดตัดสินด้วยตัวเองเสมอ
\end{itemize}

\begin{tcbinfo}{เคล็ดลับการใช้ AI}
\begin{itemize}
  \item เนื้อหาที่ AI สร้างอาจไม่แม่น (ปัญหา "ภาพหลอน") ข้อมูลสำคัญต้องตรวจสอบเสมอ
  \item อย่าป้อนข้อมูลส่วนตัวให้ AI (เลขบัตรประชาชน รหัสผ่าน ข้อมูลธนาคาร ฯลฯ)
  \item ใช้พรอมป์ (Prompt) ให้เป็น: ยิ่งอธิบายเจาะจงชัดเจนเท่าไร คุณภาพคำตอบของ AI ยิ่งสูง
  \item AI เป็นเครื่องมือ ไม่ใช่ผู้รู้แจ้ง คงการคิดและตัดสินใจด้วยตัวเองไว้
  \item เรื่องการตัดสินใจสำคัญ (การแพทย์ กฎหมาย การเงิน) ให้ปรึกษามืออาชีพ อย่าฝากไว้กับ AI ลำพัง
\end{itemize}

\end{tcbinfo}

\begin{tcbpractice}{มาลองท้าดู: โจทย์ AI ฝึกปฏิบัติสามข้อ}
\begin{enumerate}
  \item \textbf{โจทย์ที่ 1: คำสั่งหลายขั้นตอน}: พูดกับผู้ช่วยเสียงว่า: "พรุ่งนี้เช้า 7 โมงเตือนฉันให้พกร่ม และเปลี่ยนนัดในปฏิทินบ่าย 3 โมงวันนี้ไปเป็น 4 โมง" แล้วดูว่ามันทำหลายขั้นตอนในครั้งเดียวจบไหม
  \item \textbf{โจทย์ที่ 2: แต่งรูปด้วย AI}: ถ่ายภาพเที่ยวที่มีคนแปลกหน้าร่วมด้วย ใช้ฟีเจอร์ "ลบแบบมายากล" หรือ "ยางลบ AI" เก็บกวาดคนออกให้สะอาดเกลี้ยง แล้วสังเกตดูว่า AI "นึกภาพ" ฉากหลังได้เนียนแค่ไหน
  \item \textbf{โจทย์ที่ 3: ปรับแต่งพรอมป์}: ถาม AI ด้วยพรอมป์ง่าย ๆ ว่า "เขียนแผนเที่ยวหน่อย" ก่อน แล้วเขียนพรอมป์ละเอียดด้วยกรอบ CRISPE อีกครั้ง เทียบคุณภาพคำตอบของสองครั้งนี้ดู
\end{enumerate}

\end{tcbpractice}

\begin{tcbinfo}{กรณีศึกษา 2: ขั้นตอนย้ายข้อมูลเวลาเปลี่ยนเครื่องฉบับครบถ้วน}
\begin{itemize}
  \item \textbf{สถานการณ์}: เปลี่ยนจาก iPhone 13 ไปเป็น Xiaomi 15 แบบข้ามแบรนด์
  \item \textbf{ขั้นตอน}: (1) สำรองข้อมูลผ่าน iCloud/WeChat ของเครื่องเดิม (2) ออกจาก Apple ID + ปิด "ค้นหา (Find My)" (3) เครื่องใหม่เปิดแอป "ย้ายข้อมูล Xiaomi (Mi Mover)" แล้วสแกนคิวอาร์ (4) เลือกรายการที่จะย้าย (5) แชท WeChat ใช้ฟีเจอร์ "ย้ายประวัติแชท" แยกต่างหาก (6) ตรวจข้อมูลสำคัญ (รายชื่อ/อัลบั้ม/WeChat) (7) ลบข้อมูลเครื่องเดิมแบบเข้ารหัส
  \item \textbf{ผลลัพธ์ที่คาดหวัง}: ข้อมูลไม่หายสักชิ้น เครื่องเก่าขายต่อได้อย่างปลอดภัย
\end{itemize}

\end{tcbinfo}
